\chapter{Diskusija}
\label{sec:diskusija}

\TODO{
\begin{itemize}
    \item Diskusijo preusmeriti iz vprašanja, ali GNN absolutno zmaga, v vprašanje, kako se modeli obnašajo pri različnih naborih signalov.
    \item Posebej interpretirati razlike med \emph{samo pogled}, \emph{pogled in zenici} in \emph{vsi signali}: ali zenični signal pomaga nad koordinatami pogleda, ali dodatni signali v \emph{vsi signali} pomagajo ali dodajo šum, ter ali GNN bogatejšo strukturo izkoristi bolje kot negrafovski osnovni modeli.
    \item \ModelName{GazeMAE} obravnavati kot močno referenco za množico \emph{samo pogled}; če premaga modele z več signali, to interpretirati kot pomembno ugotovitev o naučenih predstavitvah pogleda.
    \item \emph{Samo zenici} in \ModelName{MOMENT} opisati samo kot prihodnje delo ali podporni poskus, če ne bosta dokončana.
\end{itemize}
}

\section{Ali grafovska predstavitev pomaga?}
\label{subsec:diskusija-graf-pomaga}

\TODO{
\begin{itemize}
    \item ali finalni GNN premaga negrafovske osnovne modele;
    \item če ne premaga: ali kaže smiselno vedenje;
    \item pri kateri nalogi je graf najbolj koristen;
    \item LOO po subjektih v primerjavi z LOO po posnetkih;
    \item primerjava proti \ModelName{GazeMAE};
    \item povezava z ablations;
    \item konservativen zaključek.
\end{itemize}
}

\section{Vloga časovnih in prostorskih povezav}
\label{subsec:diskusija-edge-types}

\TODO{
\begin{itemize}
    \item časovne povezave;
    \item prostorske povezave;
    \item fiksacijske povezave;
    \item ločevanje časovnih povezav naprej in nazaj;
    \item naučene predznačene uteži povezav;
    \item \texttt{distance-avg};
    \item \texttt{fixation-duration};
    \item značilka povezave \texttt{delta\_distance};
    \item kaj to pomeni za signale sledilnika pogleda;
    \item brez kavzalnih trditev.
\end{itemize}
}

\section{Primerjava z \ModelName{GazeMAE}}
\label{subsec:diskusija-gazemae}

\TODO{
\begin{itemize}
    \item ali \ModelName{GazeMAE} premaga negrafovske osnovne modele;
    \item ali \ModelName{GazeMAE} premaga GNN;
    \item samonadzorovano učenje predstavitev kot uporabna smer;
    \item zamrznjena enkoderja position/velocity;
    \item 512-dim okenska predstavitev;
    \item omejitev: predstavitveni model s prenosom;
    \item omejitev: ne polna reprodukcija predtreniranja.
\end{itemize}
}

\section{Stabilnost učenja in notranje predstavitve}
\label{subsec:diskusija-stabilnost}

\TODO{
\begin{itemize}
    \item Povzemi, ali krivulje izgub kažejo stabilno učenje ali prekomerno prileganje.
    \item Interpretiraj varianco in kosinusno podobnost predstavitev v povezavi s prekomernim glajenjem.
    \item Interpretiraj razpon logitov v povezavi z zanesljivostjo oziroma zasičenostjo napovedne glave.
    \item V eni povedi poročaj, da diagnostika iz dodatka~\ref{app:diagnostika-prekomerno-glajenje} kaže informativno sliko učenja: GNN se uči razločevalnih grafovskih predstavitev in pri glavnih signalnih naborih ne kaže jasnega kolapsa oziroma prekomernega glajenja.
    \item Če diagnostika ne pokaže jasnega vzorca, to napiši previdno in brez pretiranih sklepov.
    \item Poveži ugotovitve z arhitekturnimi odločitvami iz poglavja~\ref{subsec:stabilizacija-ucenja}.
\end{itemize}
}

\section{Primerjava z izvirnim \MAHNOB{} pristopom}
\label{subsec:diskusija-mahnob}

\TODO{
\begin{itemize}
    \item izvirni članek: ročno oblikovane značilke pogleda;
    \item izvirni članek: drugačen cevovod izbire značilk;
    \item izvirni članek: \ModelName{SVM};
    \item ta naloga: grafi iz časovnih oken surovejših signalov;
    \item primerjava metodološko motivacijska;
    \item ne neposredna reprodukcija;
    \item posebej omeniti uporabo \texttt{distance-avg} in fiksacijskih informacij kot približanje širšemu naboru signalov iz literature.
\end{itemize}
}

\section{Omejitve}
\label{subsec:omejitve}

Glavne omejitve naloge so:
\begin{itemize}
    \item uporaba ene glavne podatkovne zbirke;
    \item subjektivnost čustvenih oznak;
    \item omejitev na signale sledilnika pogleda brez drugih fizioloških modalnosti;
    \item občutljivost grafovske konstrukcije na hiperparametre;
    \item računska zahtevnost validacije z LOO protokoli;
    \item možnost prekomernega glajenja v grafovskih modelih;
    \item zaradi subject kfold model med treniranjem vidi patterne recordingov, tako da se lahko uči le teh. narediti bi bilo potrebno tudi CV po recordingih.
\end{itemize}

\section{Možnosti nadaljnjega dela}
\label{subsec:future-work}

Možne smeri nadaljnjega dela vključujejo:
\begin{itemize}
    \item učenje povezav namesto ročne konstrukcije grafa;
    \item uporabo grafovskih transformerjev;
    \item samonadzorovano predtreniranje na neoznačenih signalih sledilnika pogleda;
    \item razvoj grafovskega temeljnega modela za podatke pogleda;
    \item razširitev na večmodalne fiziološke podatke;
    \item evalvacijo na več podatkovnih zbirkah.
\end{itemize}

\TODO{
Pri pisanju tega odseka dodaj tudi dve konkretni smeri nadaljnjega dela:
\begin{itemize}
    \item razmislek in implementacijo uporabe mežikov, saj literatura mežike obravnava kot informativen signal pri prepoznavi čustvenih stanj iz podatkov pogleda~\cite{lim2020emotion,tarnowski2020eye};
    \item razširitev grafovske predstavitve s fiksacijskimi meta-vozlišči, torej idejo 3D lift, kjer bi posamezne meritve ostale osnovna vozlišča, fiksacije pa bi bile dodatna vozlišča višje ravni, povezana z meritvami, ki jim pripadajo.
\end{itemize}
}
